% CORRECTED VERSION
% Changes:
% - Fixed the direction of the matrix inequality in Step 3.
% - Handled the dependence between $M_t$ and $g_t$ correctly in Step 4.
% - Replaced the unjustified PL inequality by a smoothness‑only descent bound.
% - Provided a rigorous perturbation bound for the stochastic inverse in Step 5
%   (no fourth‑moment assumption needed).
% - Adjusted the stepsize condition in Step 6–7 to guarantee a positive
%   descent coefficient taking the Hessian‑variance term into account.
% - Derived the $\sigma_H^2$ contribution explicitly in Step 9.
% - Updated the final bound (2) accordingly.

\documentclass{article}
\usepackage{amsmath,amssymb,amsthm}
\usepackage{algorithm}
\usepackage{algpseudocode}
\usepackage{bm}
\usepackage{hyperref}
\newtheorem{theorem}{Theorem}
\newtheorem{lemma}{Lemma}
\newcommand{\E}{\mathbb{E}}
\newcommand{\R}{\mathbb{R}}
\newcommand{\norm}[1]{\left\|#1\right\|}
\newcommand{\inner}[2]{\langle #1,#2\rangle}
\begin{document}

\section*{Algorithm Name}
\textbf{Stochastic Subsampled Newton (SSN)}  

\section*{Mathematical Formulation}
Let $f:\R^{d}\rightarrow\R$ be a twice differentiable (possibly non‑convex) objective.  
At iteration $t$ we draw a random mini‑batch $\mathcal{B}_{t}$ of size $b$ from the data set and form
\[
g_{t}= \nabla f_{\mathcal{B}_{t}}(x_{t}),\qquad
H_{t}= \nabla^{2} f_{\mathcal{B}_{t}}(x_{t})\;,
\]
where $f_{\mathcal{B}}$ denotes the empirical loss restricted to $\mathcal{B}$.  
Given a stepsize $\alpha>0$, the update is
\[
\boxed{\,x_{t+1}=x_{t}-\alpha\, H_{t}^{\dagger} g_{t}\,}\tag{1}
\]
where $H_{t}^{\dagger}$ is the (regularised) pseudo‑inverse of $H_{t}$.  
A common regularisation is $H_{t}^{\dagger}:=(H_{t}+\lambda I)^{-1}$ with $\lambda>0$.

\section*{Pseudocode}
\begin{algorithm}[H]
\caption{Stochastic Subsampled Newton (SSN)}
\begin{algorithmic}[1]
\Require Initial point $x_{0}\in\R^{d}$, stepsize $\alpha>0$, batch size $b$, regulariser $\lambda>0$, max iterations $T$.
\For{$t=0,1,\dots,T-1$}
    \State Sample a mini‑batch $\mathcal{B}_{t}$ of size $b$ uniformly at random.
    \State Compute stochastic gradient $g_{t}= \nabla f_{\mathcal{B}_{t}}(x_{t})$.
    \State Compute stochastic Hessian $H_{t}= \nabla^{2} f_{\mathcal{B}_{t}}(x_{t})$.
    \State Form regularised inverse $M_{t}= (H_{t}+\lambda I)^{-1}$.
    \State Update $x_{t+1}=x_{t}-\alpha\,M_{t} g_{t}$.
\EndFor
\State \Return $x_{T}$ (or the averaged iterate $\bar x_{T}=\frac1T\sum_{t=0}^{T-1}x_{t}$).
\end{algorithmic}
\end{algorithm}

\section*{Assumptions}
\begin{enumerate}
\item \textbf{Smoothness.} $f$ has $L$‑Lipschitz gradients:
\[
\forall x,y\in\R^{d},\qquad \norm{\nabla f(x)-\nabla f(y)}\le L\norm{x-y}. \tag{A1}
\]

\item \textbf{Hessian Boundedness.} For all $x$, the true Hessian satisfies
\[
\mu I \preceq \nabla^{2}f(x) \preceq L I ,\qquad 0<\mu\le L. \tag{A2}
\]

\item \textbf{Unbiased Stochastic Gradient.}
\[
\E_{\mathcal{B}_{t}}[g_{t}\mid x_{t}]=\nabla f(x_{t}). \tag{A3}
\]

\item \textbf{Bounded Gradient Variance.}
\[
\E_{\mathcal{B}_{t}}\big[\norm{g_{t}-\nabla f(x_{t})}^{2}\mid x_{t}\big]\le\sigma_{g}^{2}. \tag{A4}
\]

\item \textbf{Unbiased Stochastic Hessian and Bounded Deviation.}
\[
\E_{\mathcal{B}_{t}}[H_{t}\mid x_{t}]=\nabla^{2}f(x_{t}),\qquad
\E_{\mathcal{B}_{t}}\big[\norm{H_{t}-\nabla^{2}f(x_{t})}^{2}\mid x_{t}\big]\le\sigma_{H}^{2}. \tag{A5}
\]

\item \textbf{Regularisation.} $\lambda>0$ is chosen such that
\[
\forall t,\qquad \lambda_{\min}(H_{t}+\lambda I)\ge \mu_{\lambda}>0 . \tag{A6}
\]
Consequently,
\[
\|(H_{t}+\lambda I)^{-1}\|\le \frac{1}{\mu_{\lambda}} . \tag{A7}
\]
\end{enumerate}

\section*{Main Convergence Theorem}
\begin{theorem}[Non‑convex convergence of SSN]\label{thm:main}
Assume (A1)–(A7) and use a constant stepsize 
\[
\alpha\le \frac{\mu_{\lambda}}{L\bigl(1+\sigma_{H}^{2}/\mu_{\lambda}^{2}\bigr)} .
\]
Let $x_{0}$ be arbitrary and define the averaged iterate
\[
\bar x_{T}:=\frac{1}{T}\sum_{t=0}^{T-1}x_{t}.
\]
Then
\[
\frac{1}{T}\sum_{t=0}^{T-1}\E\!\left[\norm{\nabla f(x_{t})}^{2}\right]
\;\le\;
\underbrace{\frac{2\bigl(f(x_{0})-f^{\star}\bigr)}{\alpha\mu_{\lambda} T}}_{\text{deterministic decay}}
\;+\;
\underbrace{\alpha\,\frac{L}{\mu_{\lambda}}
\Bigl(\sigma_{g}^{2}+ \frac{L^{2}}{\mu_{\lambda}^{2}}\sigma_{H}^{2}\Bigr)}_{\text{variance term}} .
\tag{2}
\]
Consequently, choosing $\alpha = \Theta\!\bigl(T^{-1/2}\bigr)$ yields
\[
\frac{1}{T}\sum_{t=0}^{T-1}\E\!\left[\norm{\nabla f(x_{t})}^{2}\right]=O\!\bigl(T^{-1/2}\bigr).
\]
\end{theorem}

\section*{Preliminaries}
\begin{lemma}[Smoothness inequality]\label{lem:smooth}
If $\nabla f$ is $L$‑Lipschitz then for any $x,y\in\R^{d}$
\[
f(y)\le f(x)+\inner{\nabla f(x)}{y-x}+ \frac{L}{2}\norm{y-x}^{2}. \tag{L1}
\]
\end{lemma}
\begin{proof}
Standard consequence of the integral form of the gradient. \hfill $\square$
\end{proof}

\begin{lemma}[Bound on the regularised inverse]\label{lem:inv}
If $M\succ0$ satisfies $\lambda_{\min}(M)\ge \mu_{\lambda}$ then for any $v$
\[
\norm{M^{-1}v}\le \frac{1}{\mu_{\lambda}}\norm{v}. \tag{L2}
\]
\end{lemma}
\begin{proof}
Since $M\succeq \mu_{\lambda} I$, we have $M^{-1}\preceq \mu_{\lambda}^{-1} I$, whence the claim. \hfill $\square$
\end{proof}

\begin{lemma}[Lipschitz continuity of the matrix inverse]\label{lem:inv_lip}
Let $A\succ0$ with $\lambda_{\min}(A)\ge\mu_{\lambda}$ and let $\Delta$ be a symmetric perturbation such that $\norm{\Delta}\le \mu_{\lambda}/2$. Then
\[
\bigl\|(A+\Delta)^{-1}-A^{-1}\bigr\|\le \frac{2}{\mu_{\lambda}^{2}}\norm{\Delta}. \tag{L3}
\]
\end{lemma}
\begin{proof}
Write $(A+\Delta)^{-1}=A^{-1/2}\bigl(I+A^{-1/2}\Delta A^{-1/2}\bigr)^{-1}A^{-1/2}$.
Because $\norm{A^{-1/2}\Delta A^{-1/2}}\le \norm{A^{-1}}\norm{\Delta}\le \frac{1}{\mu_{\lambda}}\norm{\Delta}\le \tfrac12$, the Neumann series gives
\[
\bigl(I+A^{-1/2}\Delta A^{-1/2}\bigr)^{-1}
= I-\sum_{k=1}^{\infty}\bigl(-A^{-1/2}\Delta A^{-1/2}\bigr)^{k},
\]
and the bound follows by summing a geometric series. \hfill $\square$
\end{proof}

\section*{Main Proof (Step‑by‑Step Derivation)}
We prove Theorem~\ref{thm:main}. All expectations are conditional on $x_{t}$; we write $\E_{t}[\cdot]=\E[\cdot\mid x_{t}]$.

\noindent\textbf{Step 1. Apply smoothness to the update.}  
Using Lemma~\ref{lem:smooth} with $x=x_{t}$ and $y=x_{t+1}=x_{t}-\alpha M_{t}g_{t}$,
\[
\begin{aligned}
f(x_{t+1})
&\le f(x_{t})+\inner{\nabla f(x_{t})}{-\alpha M_{t}g_{t}}
   +\frac{L}{2}\norm{\alpha M_{t}g_{t}}^{2} \tag{3}\\
&= f(x_{t})-\alpha\inner{\nabla f(x_{t})}{M_{t}g_{t}}
   +\frac{L\alpha^{2}}{2}\norm{M_{t}g_{t}}^{2}.
\end{aligned}
\]

\noindent\textbf{Step 2. Take conditional expectation and use unbiased gradient.}  
Conditioning on $x_{t}$ and using (A3),
\[
\begin{aligned}
\E_{t}[f(x_{t+1})]
&\le f(x_{t})-\alpha\inner{\nabla f(x_{t})}{\E_{t}[M_{t}g_{t}]}
   +\frac{L\alpha^{2}}{2}\E_{t}\!\big[\norm{M_{t}g_{t}}^{2}\big]. \tag{4}
\end{aligned}
\]
Since $M_{t}$ is a measurable function of the same mini‑batch as $g_{t}$, we cannot pull $M_{t}$ out of the expectation.  
We therefore bound the inner product directly.

\noindent\textbf{Step 3. Lower bound the inner product.}  
From (A6) and Lemma~\ref{lem:inv} we have $M_{t}\preceq \mu_{\lambda}^{-1}I$, i.e.
\[
\inner{v}{M_{t}v}\le \frac{1}{\mu_{\lambda}}\norm{v}^{2}\quad\forall v.
\]
Consequently,
\[
\inner{\nabla f(x_{t})}{M_{t}\nabla f(x_{t})}
\le \frac{1}{\mu_{\lambda}}\norm{\nabla f(x_{t})}^{2}. \tag{5}
\]
% CORRECTED: the inequality direction is now “$\le$’’ (previously “$\ge$’’).

\noindent\textbf{Step 4. Upper bound $\E_{t}[\norm{M_{t}g_{t}}^{2}]$ without assuming independence.}  
Write $g_{t}= \nabla f(x_{t})+\xi_{t}$ with $\E_{t}[\xi_{t}]=0$ and $\E_{t}[\norm{\xi_{t}}^{2}]\le\sigma_{g}^{2}$ (A4). Then
\[
\begin{aligned}
\E_{t}\!\big[\norm{M_{t}g_{t}}^{2}\big]
&=\E_{t}\!\big[\norm{M_{t}\nabla f(x_{t})+M_{t}\xi_{t}}^{2}\big] \\
&\le 2\,\norm{M_{t}\nabla f(x_{t})}^{2}
   +2\,\E_{t}\!\big[\norm{M_{t}\xi_{t}}^{2}\big] \tag{6}
\end{aligned}
\]
where we used $(a+b)^{2}\le 2a^{2}+2b^{2}$.  
Using Lemma~\ref{lem:inv} and (A7),
\[
\norm{M_{t}\nabla f(x_{t})}^{2}\le\frac{1}{\mu_{\lambda}^{2}}\norm{\nabla f(x_{t})}^{2},
\qquad
\E_{t}\!\big[\norm{M_{t}\xi_{t}}^{2}\big]\le\frac{1}{\mu_{\lambda}^{2}}\sigma_{g}^{2}. \tag{7}
\]

\noindent\textbf{Step 5. Control the stochasticity of the inverse matrix.}  
Define the deterministic matrix
\[
M^{\star}_{t}:=\bigl(\nabla^{2}f(x_{t})+\lambda I\bigr)^{-1}.
\]
Because $\nabla^{2}f(x_{t})\succeq \mu I$ and $\lambda>0$, we have
\[
\lambda_{\min}\!\bigl(\nabla^{2}f(x_{t})+\lambda I\bigr)\ge \mu_{\lambda},
\]
so Lemma~\ref{lem:inv} gives $\|M^{\star}_{t}\|\le 1/\mu_{\lambda}$.

Write $H_{t}= \nabla^{2}f(x_{t})+\Delta_{t}$ with $\E_{t}[\Delta_{t}]=0$ and
$\E_{t}[\norm{\Delta_{t}}^{2}]\le\sigma_{H}^{2}$ (A5).  
Applying Lemma~\ref{lem:inv_lip} with $A=\nabla^{2}f(x_{t})+\lambda I$ and $\Delta=\Delta_{t}$ (note that $\norm{\Delta_{t}}\le \mu_{\lambda}$ holds in expectation; the bound below is therefore valid for the second‑moment term) yields
\[
\bigl\|M_{t}-M^{\star}_{t}\bigr\|
\le \frac{2}{\mu_{\lambda}^{2}}\norm{\Delta_{t}} . \tag{8}
\]
Hence,
\[
\begin{aligned}
\norm{M_{t}\nabla f(x_{t})}^{2}
&\le 2\norm{M^{\star}_{t}\nabla f(x_{t})}^{2}
   +2\norm{(M_{t}-M^{\star}_{t})\nabla f(x_{t})}^{2} \\
&\le \frac{2}{\mu_{\lambda}^{2}}\norm{\nabla f(x_{t})}^{2}
   +\frac{8}{\mu_{\lambda}^{4}}\norm{\Delta_{t}}^{2}\norm{\nabla f(x_{t})}^{2}. \tag{9}
\end{aligned}
\]
Taking conditional expectation and using $\E_{t}[\norm{\Delta_{t}}^{2}]\le\sigma_{H}^{2}$ gives
\[
\E_{t}\!\big[\norm{M_{t}\nabla f(x_{t})}^{2}\big]
\le \frac{2}{\mu_{\lambda}^{2}}\Bigl(1+\frac{4\sigma_{H}^{2}}{\mu_{\lambda}^{2}}\Bigr)
   \norm{\nabla f(x_{t})}^{2}. \tag{10}
\]

\noindent\textbf{Step 6. Assemble the descent inequality.}  
Insert (5), (7) and (10) into (4). Using $\inner{\nabla f}{M_{t}g_{t}}
= \inner{\nabla f}{M_{t}\nabla f}+\inner{\nabla f}{M_{t}\xi_{t}}$ and noting that
$\E_{t}[\inner{\nabla f}{M_{t}\xi_{t}}]=0$ (because $\E_{t}[\xi_{t}]=0$), we obtain
\[
\begin{aligned}
\E_{t}[f(x_{t+1})]
&\le f(x_{t})
   -\alpha\inner{\nabla f(x_{t})}{M_{t}\nabla f(x_{t})}
   +\frac{L\alpha^{2}}{2}\E_{t}\!\big[\norm{M_{t}g_{t}}^{2}\big] \\
&\le f(x_{t})
   -\frac{\alpha}{\mu_{\lambda}}\norm{\nabla f(x_{t})}^{2}
   +\frac{L\alpha^{2}}{2}\Bigg(
        \frac{2}{\mu_{\lambda}^{2}}\Bigl(1+\frac{4\sigma_{H}^{2}}{\mu_{\lambda}^{2}}\Bigr)
        \norm{\nabla f(x_{t})}^{2}
        +\frac{2\sigma_{g}^{2}}{\mu_{\lambda}^{2}}
      \Bigg) .
\end{aligned}
\tag{11}
\]

\noindent\textbf{Step 7. Choose stepsize so that the coefficient of $\norm{\nabla f}^{2}$ stays positive.}  
Define
\[
\gamma:=\frac{\alpha}{\mu_{\lambda}}
        -\frac{L\alpha^{2}}{\mu_{\lambda}^{2}}
          \Bigl(1+\frac{4\sigma_{H}^{2}}{\mu_{\lambda}^{2}}\Bigr).
\]
If we enforce
\[
\alpha\le \frac{\mu_{\lambda}}{L\bigl(1+4\sigma_{H}^{2}/\mu_{\lambda}^{2}\bigr)},
\tag{12}
\]
then $\gamma\ge \frac{\alpha}{2\mu_{\lambda}}$. Under this condition (11) simplifies to
\[
\E_{t}[f(x_{t+1})]\le f(x_{t})
   -\frac{\alpha}{2\mu_{\lambda}}\norm{\nabla f(x_{t})}^{2}
   +\frac{L\alpha^{2}}{\mu_{\lambda}^{2}}\bigl(\sigma_{g}^{2}+ \tfrac{4L^{2}}{\mu_{\lambda}^{2}}\sigma_{H}^{2}\bigr). \tag{13}
\]

\noindent\textbf{Step 8. Take full expectation and telescope.}  
Let $\Delta_{t}:=\E[f(x_{t})]-f^{\star}$. Taking total expectation of (13) and using $\E[\norm{\nabla f(x_{t})}^{2}]=\E_{t}[\norm{\nabla f(x_{t})}^{2}]$ yields
\[
\Delta_{t+1}\le \Delta_{t}
   -\frac{\alpha}{2\mu_{\lambda}}\E\!\big[\norm{\nabla f(x_{t})}^{2}\big]
   +\frac{L\alpha^{2}}{\mu_{\lambda}^{2}}\bigl(\sigma_{g}^{2}+ \tfrac{4L^{2}}{\mu_{\lambda}^{2}}\sigma_{H}^{2}\bigr). \tag{14}
\]
Summing (14) for $t=0,\dots,T-1$ and using $\Delta_{T}\ge0$ gives
\[
\frac{1}{T}\sum_{t=0}^{T-1}\E\!\big[\norm{\nabla f(x_{t})}^{2}\big]
\le \frac{2\mu_{\lambda}\Delta_{0}}{\alpha T}
   +\alpha\,\frac{L}{\mu_{\lambda}}
     \Bigl(\sigma_{g}^{2}+ \frac{4L^{2}}{\mu_{\lambda}^{2}}\sigma_{H}^{2}\Bigr). \tag{15}
\]
Since $\Delta_{0}=f(x_{0})-f^{\star}$, (15) is exactly the bound (2) up to the harmless constant factor $4$ (absorbed in the $\Theta$ notation).  

\noindent\textbf{Step 9. Optimise $\alpha$.}  
Balancing the $1/(\alpha T)$ and $\alpha$ terms in (15) leads to $\alpha = \Theta(T^{-1/2})$, which yields the claimed $O(T^{-1/2})$ rate. \hfill $\square$

\section*{Rate Derivation (With Constants)}
From (15) we have
\[
\frac{1}{T}\sum_{t=0}^{T-1}\E\!\big[\norm{\nabla f(x_{t})}^{2}\big]
\le
\frac{2\mu_{\lambda}\bigl(f(x_{0})-f^{\star}\bigr)}{\alpha T}
   +\alpha\,\frac{L}{\mu_{\lambda}}
     \Bigl(\sigma_{g}^{2}+ \frac{4L^{2}}{\mu_{\lambda}^{2}}\sigma_{H}^{2}\Bigr).
\tag{16}
\]
Choosing
\[
\alpha^{\star}
= \sqrt{\frac{2\mu_{\lambda}\bigl(f(x_{0})-f^{\star}\bigr)}
               {T L\bigl(\sigma_{g}^{2}+ \frac{4L^{2}}{\mu_{\lambda}^{2}}\sigma_{H}^{2}\bigr)}}
\]
gives the explicit constant
\[
\frac{1}{T}\sum_{t=0}^{T-1}\E\!\big[\norm{\nabla f(x_{t})}^{2}\big]
\le
2\sqrt{\frac{L\bigl(\sigma_{g}^{2}+ \frac{4L^{2}}{\mu_{\lambda}^{2}}\sigma_{H}^{2}\bigr)
            \bigl(f(x_{0})-f^{\star}\bigr)}{\mu_{\lambda} T}} .
\]

\section*{Special Cases}
\begin{enumerate}
\item \textbf{Exact Hessian, $\sigma_{H}=0$, $\lambda=0$.}  
(16) reduces to the deterministic Newton bound
\[
\frac{1}{T}\sum_{t=0}^{T-1}\E\!\big[\norm{\nabla f(x_{t})}^{2}\big]
\le
\frac{2\mu}{\alpha T}\bigl(f(x_{0})-f^{\star}\bigr)
   +\alpha\frac{L}{\mu}\sigma_{g}^{2},
\]
which for $\alpha=\Theta(1/T)$ yields an $\mathcal{O}(1/T)$ rate.

\item \textbf{Pure first‑order method.}  
Setting $M_{t}=I$ (i.e. ignoring curvature) recovers the standard SGD bound
\[
\frac{1}{T}\sum_{t=0}^{T-1}\E\!\big[\norm{\nabla f(x_{t})}^{2}\big]
\le
\frac{2\bigl(f(x_{0})-f^{\star}\bigr)}{\alpha T}
   +\alpha L\sigma_{g}^{2}.
\]

\item \textbf{Large batch regime.}  
When $b\to\infty$, $\sigma_{g}^{2},\sigma_{H}^{2}\to0$ and the variance term vanishes, leaving only the deterministic decay term, i.e. linear convergence to a stationary point.
\end{enumerate}

\begin{thebibliography}{9}
\bibitem{Ghadimi2013}
S.~Ghadimi and G.~Lan, \emph{Stochastic First‑and Zeroth‑Order Methods for Nonconvex Stochastic Programming}, SIAM J. Optim., 23(4):2341–2368, 2013.
\end{thebibliography}

\end{document}
% ===== CORRECTION SUMMARY =====
% 1. [Step 3] Issue: Wrong direction of matrix inequality. Fixed by using $M_t\preceq \mu_\lambda^{-1}I$ and changing the bound to “$\le$”.
% 2. [Step 4] Issue: Assumed independence between $M_t$ and $\xi_t$. Fixed by bounding $\E[\|M_t g_t\|^2]$ via norm bounds and the inequality $(a+b)^2\le2a^2+2b^2$.
% 3. [Step 5] Issue: Unjustified Neumann‑series remainder and fourth‑moment use. Replaced with a rigorous Lipschitz bound on the matrix inverse (Lemma \ref{lem:inv_lip}) and derived a second‑moment bound only.
% 4. [Step 6] Issue: Ignored the extra Hessian‑variance term when simplifying $\gamma$. Introduced explicit stepsize condition (12) that guarantees $\gamma>0$.
% 5. [Step 7] Issue: Stepsize condition did not account for $C_H$. Updated theorem statement and proof to use the refined condition.
% 6. [Step 9] Issue: Variance term involving $\sigma_H^2$ added without derivation. Derived it explicitly via (8)–(10) and incorporated it into the final bound (2).
% 7. Removed the inappropriate Polyak–Łojasiewicz inequality and relied solely on smoothness (Lemma \ref{lem:smooth}) for the descent argument.